\section{Failure Analysis and Design Implications}
\label{sec:failure}

\subsection{Route Viability Precedes Router Optimization}

The router failure was not simply a poor choice of $\tau$. Validation success was
8.75\% at every candidate threshold from 0.60 through 1.00. One deployable family
had no evidence-task success; the other produced a repeated, unparsable sequence
under the frozen integration. No calibration method can create a useful operating
point when every feasible action is weak or operationally incompatible. A practical
router-development pipeline should therefore impose a precondition: at least two
deployable actions must demonstrate nontrivial, query-dependent strict success
before policy fitting begins. If the condition fails, the correct output is an
infeasible action pool, not a forced router leaderboard entry.

The exact-matrix evaluation also clarifies what happened on test. The policy did
not prospectively discover that A03 was optimal for every question; it invoked the
same safe fallback because no predicted action cleared the fail-closed threshold.
This distinction matters when diagnosing routing systems and when comparing them
with simple baselines \cite{LiEtAl2026LLMRouterBench}.

\subsection{Grounding Is a Generator--Task Interaction}

The aggregate transfer effects conceal sharp task interactions. Always grounding
helped the 3B/4B generators on several evidence-requiring classes while removing
some direct successes on no-retrieval and deleted/missing cases. For the 0.6B model,
it removed the only successful stratum and failed to unlock any other one. An
adaptive system therefore needs both a capable grounded route and a reliable
retrieval-need signal. ``Small model'' is too coarse an explanatory variable;
prompt compliance, context use, instruction tuning, tokenizer/chat templates, and
task structure jointly determine whether evidence becomes useful.

This finding also limits the common practice of reporting one RAG-minus-direct
average. Such an average depends on the task mixture. Under traffic dominated by
knowledge or temporal questions, the larger grounded endpoints would look more
favorable; under direct-copy or valid-abstention traffic, always-grounded routes
would look worse. Deployment claims require either a representative workload or
stratified reporting.

\subsection{The Interface Is Part of the System}

Strict success exposed output-interface failures that answer-only evaluation would
blur. The Granite-4.0-1B repetition failure occurred before semantic evaluation.
A03 often produced superficially citation-like values that did not match exact
retrieved IDs. The deterministic guard then rejected unsupported or malformed
outputs. These are not cosmetic parser errors: a production evidence system needs
machine-actionable citations and predictable abstention. Still, exact-ID scoring
can be harsher than human recoverability, so future work should separate schema
conformance from semantic citation entailment instead of treating either as the
other.

The observed failure suggests a mandatory compatibility gate per
model--prompt--runtime triple: validate the chat template, schema adherence,
termination behavior, and citation serialization before collecting expensive
route matrices. A compatibility failure should create a new versioned experimental
specification rather than a silent prompt repair after test access.

\subsection{Measured Energy Is Not Total RAG Energy}

Every grounded endpoint increased the energy observed during generation, consistent
with longer prompts and, frequently, longer outputs. This does not imply that the
reported difference equals the total energy cost of RAG. CPU dense encoding,
lexical search, memory traversal, routing, and verification are absent from the
energy quantity. Nor are host idle power, storage, cooling, and network transfers.
The result supports a narrower statement: under these frozen routes and same-board
matched pairs, generation itself consumed more selected-board energy when the
grounded prompt was used.

Future evaluations should add synchronized CPU/host power measurement, idle-power
subtraction where appropriate, repeated interleaved trials, and an explicit
boundary diagram. Reporting per-query energy together with token counts, latency,
and useful successes would make overhead and output-length effects more legible.

\subsection{Robustness Requires a Nonzero Clean Reference}

A zero corruption delta is meaningful only when the clean system performs the
task. Here the selected route began at zero, so the corruption study cannot
distinguish stable failure from robust competence. A better robustness gate would
first require a minimum clean utility for each evaluated anchor, then report both
absolute corrupted success and paired degradation. Attack-specific compromise
metrics should likewise be presented as narrow detectors, not substitutes for
task success or general adversarial robustness.

\subsection{Actionable Recommendations}

The evidence supports five concrete design rules:

\begin{enumerate}
    \item gate model/chat-template/schema compatibility before route benchmarking;
    \item establish generator- and task-specific grounding value before training a
    controller;
    \item treat validation infeasibility as a first-class result with a fail-closed
    deployment state;
    \item measure and label energy boundaries precisely, and avoid causal energy
    comparisons across unbalanced boards or sessions;
    \item require nonzero clean utility before interpreting robustness deltas.
\end{enumerate}
